\chapter{Complexity — CNF Encoding}
%%% generated by the blueprint pipeline from TCSlib.Complexity.Formulas.CNFEncoding
%%%%%%%%%%%%%%%%%%%%%%%%%%%%%%%%%%%%%%%%%%%%%%%%%%%%%%%%%%%%%%%%%%%%%%%%%%%%%%%
\section{Overview}

Since $\mathrm{SAT}$ and $3\mathrm{SAT}$ are languages of bit strings [AB09, Section~2.3.1],
this module fixes a concrete binary encoding of CNF formulas over variables indexed by
natural numbers: a literal $(v,b)$ is written as $1^{v+1}0b$ (the variable index in unary),
a clause as its literal encodings followed by $0$, and a formula as its clause encodings,
each prefixed by $1$, followed by $0$. It provides an exact-consumption parser for this
grammar and a total decoder sending ill-formed strings to the empty formula
[AB09, footnote~3], and proves that decoding inverts encoding and that the formula decoded
from a string $x$ has variable count at most $|x|$.

Throughout, bit strings are finite words over $\{0,1\}$ (with $1$ standing for true and $0$
for false), a literal is a pair $(v,b)$ of a variable index $v \in \mathbb{N}$ and a polarity
bit $b \in \{0,1\}$, a clause is a finite list of literals, and a CNF formula is a finite
list of clauses.

%%%%%%%%%%%%%%%%%%%%%%%%%%%%%%%%%%%%%%%%%%%%%%%%%%%%%%%%%%%%%%%%%%%%%%%%%%%%%%%
\section{Declarations}

\begin{definition}[Binary encoding of a literal]
\lean{Std.Sat.CNF.serializeLit}
\label{Std.Sat.CNF.serializeLit}
\leanok
For a literal $\ell = (v, b)$, consisting of a variable index $v \in \mathbb{N}$ and a
polarity bit $b \in \{0,1\}$, its binary encoding is the bit string
$\mathrm{enc}(\ell) = 1^{v+1}\,0\,b$ of length $v+3$ (writing $1$ for true and $0$ for
false): the index in unary as $v+1$ ones, nonempty even when $v = 0$, then a terminating
$0$, then the polarity bit $b$.
\end{definition}

\begin{definition}[Binary encoding of a clause]
\lean{Std.Sat.CNF.serializeClause}
\label{Std.Sat.CNF.serializeClause}
\leanok
\uses{Std.Sat.CNF.serializeLit}
Let $C = (\ell_1, \dots, \ell_m)$ be a clause, i.e.\ a finite list of literals
$\ell_i = (v_i, b_i)$ with variable indices $v_i \in \mathbb{N}$ and polarity bits
$b_i \in \{0,1\}$. Its binary encoding is the bit string
\[
  \mathrm{enc}(C) = 1^{v_1+1}0\,b_1\;1^{v_2+1}0\,b_2 \cdots 1^{v_m+1}0\,b_m\;0,
\]
the concatenation of the literal encodings $1^{v_i+1}0b_i$ followed by a closing $0$; the
closing bit is unambiguous because every literal encoding begins with $1$.
\end{definition}

\begin{definition}[Binary encoding of a CNF formula]
\lean{Std.Sat.CNF.serialize}
\label{Std.Sat.CNF.serialize}
\leanok
\uses{Std.Sat.CNF.serializeClause}
Let $\varphi = (C_1, \dots, C_k)$ be a CNF formula over variables indexed by $\mathbb{N}$,
given as a finite list of clauses, each a finite list of literals $(v,b)$ with
$v \in \mathbb{N}$ and $b \in \{0,1\}$. Its binary encoding is
\[
  \mathrm{enc}(\varphi) = 1\,\mathrm{enc}(C_1)\;1\,\mathrm{enc}(C_2) \cdots 1\,\mathrm{enc}(C_k)\;0,
\]
where a clause $(\ell_1,\dots,\ell_m)$ is encoded as
$\mathrm{enc}(\ell_1)\cdots\mathrm{enc}(\ell_m)\,0$ and a literal $(v,b)$ as $1^{v+1}0b$.
At the formula level a $1$ announces another clause and a $0$ ends the formula; in
particular the empty formula is encoded as the single bit $0$.
\end{definition}

\begin{definition}[Leading run of ones]
\lean{Std.Sat.CNF.takeTrues}
\label{Std.Sat.CNF.takeTrues}
\leanok
For a bit string $x$, this operation returns the pair $(k, y)$ where $k$ is the length of
the longest prefix of $x$ consisting only of $1$s and $y$ is the rest of $x$, so that
$x = 1^k y$ and $y$ is either empty or begins with $0$.
\end{definition}

\begin{definition}[Parsing one literal]
\lean{Std.Sat.CNF.parseLit}
\label{Std.Sat.CNF.parseLit}
\leanok
\uses{Std.Sat.CNF.takeTrues}
For a bit string $x$, the literal parser succeeds exactly when $x = 1^{k+1}\,0\,b\,y$ for
some $k \in \mathbb{N}$, bit $b \in \{0,1\}$ and bit string $y$, in which case it returns
the literal $(k, b)$, with variable index $k$ and polarity $b$, together with the
unconsumed remainder $y$. On every other input (one not beginning with $1$, or one ending
before the $0$ and the polarity bit that follow its leading run of ones) it fails.
\end{definition}

\begin{definition}[Fuel-bounded parsing of one clause]
\lean{Std.Sat.CNF.parseClause}
\label{Std.Sat.CNF.parseClause}
\leanok
\uses{Std.Sat.CNF.parseLit}
For a fuel bound $f \in \mathbb{N}$ and a bit string $x$, the clause parser with fuel $f$ is
defined by recursion on $f$. If $x = 0y$, it succeeds with the empty clause and remainder
$y$. If $f = g + 1$ and $x$ begins with $1$, it reads one literal from the front of $x$,
which succeeds only when $x = 1^{k+1}0by$ and then gives the literal $(k,b)$ and remainder
$y$; it then runs the clause parser with fuel $g$ on $y$, and if that returns a clause $C$
with remainder $z$, the result is the clause obtained by prepending $(k,b)$ to $C$, with
remainder $z$. In every other case (empty input, fuel $0$ on an input beginning with $1$,
or a failed sub-parse) it fails.
\end{definition}

\begin{definition}[Fuel-bounded parsing of a clause list]
\lean{Std.Sat.CNF.parseClauses}
\label{Std.Sat.CNF.parseClauses}
\leanok
\uses{Std.Sat.CNF.parseClause}
For a fuel bound $f \in \mathbb{N}$ and a bit string $x$, the formula parser with fuel $f$
is defined by recursion on $f$. If $x = 0y$, it succeeds with the empty formula (no
clauses) and remainder $y$. If $f = g+1$ and $x = 1y$, it runs the clause parser with fuel
$g$ on $y$ and, if this returns a clause $C$ with remainder $z$, runs the formula parser
with fuel $g$ on $z$; if that returns a formula $\psi$ with remainder $w$, the result is
the formula whose first clause is $C$ followed by the clauses of $\psi$, with remainder
$w$. In all other cases it fails. Here the clause parser with fuel $g$ on a bit string $u$
returns the empty clause and remainder $u'$ when $u = 0u'$; when $g = h+1$ and
$u = 1^{j+1}0cu'$ it prepends the literal $(j,c)$ to the clause returned by the clause
parser with fuel $h$ on $u'$, keeping that call's remainder; and it fails otherwise.
\end{definition}

\begin{definition}[Exact-consumption parser for CNF formulas]
\lean{Std.Sat.CNF.parse}
\label{Std.Sat.CNF.parse}
\leanok
\uses{Std.Sat.CNF.parseClauses}
For a bit string $x$, the parser runs the fuel-bounded formula parser with fuel $|x|$ on $x$
and returns the resulting CNF formula $\varphi$ if that run succeeds with empty remainder,
i.e.\ consumes all of $x$; it fails otherwise, in particular whenever trailing bits remain.
The formula parser reads the grammar
$\langle\text{formula}\rangle ::= 0 \mid 1\,\langle\text{clause}\rangle\,\langle\text{formula}\rangle$,
$\langle\text{clause}\rangle ::= 0 \mid \langle\text{literal}\rangle\,\langle\text{clause}\rangle$,
$\langle\text{literal}\rangle ::= 1^{k+1}\,0\,b$ (read as the literal with variable index
$k \in \mathbb{N}$ and polarity $b \in \{0,1\}$), where the fuel bounds the recursion: with
fuel $g+1$, after reading the leading $1$ of a clause the formula parser parses that clause
and the remaining formula each with fuel $g$, and after reading a literal the clause parser
parses the rest of the clause with fuel $g$; with fuel $0$ only a terminating $0$ can be
read.
\end{definition}

\begin{definition}[Fallback formula]
\lean{Std.Sat.CNF.fallback}
\label{Std.Sat.CNF.fallback}
\leanok
The fallback CNF formula over variables indexed by $\mathbb{N}$ is the empty formula,
consisting of no clauses; it is therefore satisfied by every assignment. It is the fixed
formula to which bit strings that fail to parse are decoded, following [AB09, footnote~3].
\end{definition}

\begin{definition}[Total decoding of bit strings into CNF formulas]
\lean{Std.Sat.CNF.decode}
\label{Std.Sat.CNF.decode}
\leanok
\uses{Std.Sat.CNF.fallback, Std.Sat.CNF.parse}
For a bit string $x$, its decoding $\mathrm{dec}(x)$ is the CNF formula returned by the
exact-consumption parser on $x$ when that parse succeeds, and the empty formula (no clauses)
otherwise. Here the parser reads $x$, with recursion fuel $|x|$, according to the grammar
$\langle\text{formula}\rangle ::= 0 \mid 1\,\langle\text{clause}\rangle\,\langle\text{formula}\rangle$,
$\langle\text{clause}\rangle ::= 0 \mid \langle\text{literal}\rangle\,\langle\text{clause}\rangle$,
$\langle\text{literal}\rangle ::= 1^{k+1}\,0\,b$ (the literal with variable index
$k \in \mathbb{N}$ and polarity $b \in \{0,1\}$), and succeeds only if it consumes all of
$x$.
\end{definition}

\begin{theorem}[Leading run of a unary block]
\lean{Std.Sat.CNF.takeTrues_replicate}
\label{Std.Sat.CNF.takeTrues_replicate}
\leanok
\uses{Std.Sat.CNF.takeTrues}
\difficulty{3}
For every $k \in \mathbb{N}$ and every bit string $r$, splitting the bit string
$1^{k+1}\,0\,r$ into its maximal leading run of $1$s and the remainder gives the run length
$k+1$ and the remainder $0\,r$.
\end{theorem}

\begin{theorem}[Literal round trip]
\lean{Std.Sat.CNF.parseLit_serializeLit}
\label{Std.Sat.CNF.parseLit_serializeLit}
\leanok
\uses{Std.Sat.CNF.parseLit, Std.Sat.CNF.serializeLit, Std.Sat.CNF.takeTrues_replicate}
\difficulty{1}
Let $\ell = (v, b)$ be a literal with variable index $v \in \mathbb{N}$ and polarity bit
$b \in \{0,1\}$, whose binary encoding is $1^{v+1}0b$, and let $r$ be any bit string. The
literal parser, which on an input of the form $1^{k+1}0cy$ returns the literal $(k,c)$ with
remainder $y$ and fails on every other input, succeeds on $1^{v+1}0b\,r$ and returns $\ell$
with remainder $r$.
\end{theorem}

\begin{theorem}[Clause round trip with suffix and fuel]
\lean{Std.Sat.CNF.parseClause_serializeClause}
\label{Std.Sat.CNF.parseClause_serializeClause}
\leanok
\uses{Std.Sat.CNF.parseClause, Std.Sat.CNF.parseLit_serializeLit, Std.Sat.CNF.serializeClause, Std.Sat.CNF.serializeLit}
\difficulty{4}
Let $C = (\ell_1, \dots, \ell_m)$ be a clause of literals $\ell_i = (v_i, b_i)$ with
$v_i \in \mathbb{N}$ and $b_i \in \{0,1\}$, and let
$\mathrm{enc}(C) = 1^{v_1+1}0b_1\cdots1^{v_m+1}0b_m\,0$ be its binary encoding. For every
bit string $r$ and every fuel $f \ge |\mathrm{enc}(C)|$, the clause parser with fuel $f$
succeeds on $\mathrm{enc}(C)\,r$ and returns $C$ with remainder $r$. Here the clause parser
with fuel $f$ on an input $x$ returns the empty clause and remainder $y$ when $x = 0y$; when
$f = g+1$ and $x = 1^{k+1}0cy$ it prepends the literal $(k,c)$ to the clause returned by the
clause parser with fuel $g$ on $y$, keeping that call's remainder and failing if it fails;
and it fails on every other input.
\end{theorem}

\begin{theorem}[Formula round trip with suffix and fuel]
\lean{Std.Sat.CNF.parseClauses_serialize}
\label{Std.Sat.CNF.parseClauses_serialize}
\leanok
\uses{Std.Sat.CNF.parseClause_serializeClause, Std.Sat.CNF.parseClauses, Std.Sat.CNF.serialize, Std.Sat.CNF.serializeClause}
\difficulty{4}
Let $\varphi = (C_1,\dots,C_k)$ be a CNF formula over variables indexed by $\mathbb{N}$, and
let $\mathrm{enc}(\varphi) = 1\,\mathrm{enc}(C_1)\cdots1\,\mathrm{enc}(C_k)\,0$ be its binary
encoding, where a clause with literals $(v_1,b_1),\dots,(v_m,b_m)$ is encoded as
$1^{v_1+1}0b_1\cdots1^{v_m+1}0b_m\,0$. For every bit string $r$ and every fuel
$f \ge |\mathrm{enc}(\varphi)|$, the formula parser with fuel $f$ succeeds on
$\mathrm{enc}(\varphi)\,r$ and returns $\varphi$ with remainder $r$. Here the formula parser
with fuel $f$ on an input $x$ returns the empty formula and remainder $y$ when $x = 0y$;
when $f = g+1$ and $x = 1y$ it parses a clause $C$ from $y$ with fuel $g$, leaving $z$, then
a formula $\psi$ from $z$ with fuel $g$, leaving $w$, and returns $C$ followed by the clauses
of $\psi$ with remainder $w$; and it fails otherwise. The clause parser with fuel $g$ on $u$
returns the empty clause and remainder $u'$ when $u = 0u'$, and when $g = h+1$ and
$u = 1^{j+1}0cu'$ it prepends the literal $(j,c)$ to the clause parsed from $u'$ with fuel
$h$, keeping that remainder; it fails otherwise.
\end{theorem}

\begin{theorem}[Parsing inverts the binary encoding]
\lean{Std.Sat.CNF.parse_serialize}
\label{Std.Sat.CNF.parse_serialize}
\leanok
\uses{Std.Sat.CNF.parse, Std.Sat.CNF.parseClauses_serialize, Std.Sat.CNF.serialize}
\difficulty{3}
For every CNF formula $\varphi$ over variables indexed by $\mathbb{N}$, the exact-consumption
parser applied to the binary encoding $\mathrm{enc}(\varphi)$ succeeds and returns
$\varphi$. Here a literal $(v,b)$ is encoded as $1^{v+1}0b$, a clause as the concatenation
of its literal encodings followed by $0$, and a formula $(C_1,\dots,C_k)$ as
$1\,\mathrm{enc}(C_1)\cdots1\,\mathrm{enc}(C_k)\,0$; the parser on a bit string $x$ reads
$x$, with recursion fuel $|x|$, according to the grammar
$\langle\text{formula}\rangle ::= 0 \mid 1\,\langle\text{clause}\rangle\,\langle\text{formula}\rangle$,
$\langle\text{clause}\rangle ::= 0 \mid \langle\text{literal}\rangle\,\langle\text{clause}\rangle$,
$\langle\text{literal}\rangle ::= 1^{k+1}\,0\,b$ (the literal $(k,b)$), and succeeds only if
it consumes all of $x$.
\end{theorem}

\begin{theorem}[Decoding inverts the binary encoding]
\lean{Std.Sat.CNF.decode_serialize}
\label{Std.Sat.CNF.decode_serialize}
\leanok
\uses{Std.Sat.CNF.decode, Std.Sat.CNF.parse_serialize, Std.Sat.CNF.serialize}
\difficulty{1}
For every CNF formula $\varphi$ over variables indexed by $\mathbb{N}$,
$\mathrm{dec}(\mathrm{enc}(\varphi)) = \varphi$. Here a literal $(v,b)$ is encoded as
$1^{v+1}0b$, a clause as the concatenation of its literal encodings followed by $0$, and a
formula $(C_1,\dots,C_k)$ as $1\,\mathrm{enc}(C_1)\cdots1\,\mathrm{enc}(C_k)\,0$; and
$\mathrm{dec}(x)$ is the formula obtained by parsing the bit string $x$, with recursion fuel
$|x|$, according to the grammar
$\langle\text{formula}\rangle ::= 0 \mid 1\,\langle\text{clause}\rangle\,\langle\text{formula}\rangle$,
$\langle\text{clause}\rangle ::= 0 \mid \langle\text{literal}\rangle\,\langle\text{clause}\rangle$,
$\langle\text{literal}\rangle ::= 1^{k+1}\,0\,b$ (the literal $(k,b)$) when this parse
consumes all of $x$, and the empty formula otherwise.
\end{theorem}

\begin{theorem}[Leading-run split preserves length]
\lean{Std.Sat.CNF.takeTrues_length}
\label{Std.Sat.CNF.takeTrues_length}
\leanok
\uses{Std.Sat.CNF.takeTrues}
\difficulty{4}
For every bit string $x$, if $k$ is the length of the longest prefix of $x$ consisting only
of $1$s and $y$ is the remainder of $x$ after that prefix, then $k + |y| = |x|$.
\end{theorem}

\begin{theorem}[Length consumed by a literal parse]
\lean{Std.Sat.CNF.parseLit_length}
\label{Std.Sat.CNF.parseLit_length}
\leanok
\uses{Std.Sat.CNF.parseLit, Std.Sat.CNF.takeTrues_length}
\difficulty{4}
Let $x$ and $r$ be bit strings and let $(v, b)$ be a literal with $v \in \mathbb{N}$ and
$b \in \{0,1\}$. If the literal parser, which on an input $1^{k+1}0cy$ returns the literal
$(k,c)$ with remainder $y$ and fails on every other input, returns $(v,b)$ with remainder
$r$ when applied to $x$, then $v + 3 + |r| = |x|$.
\end{theorem}

\begin{theorem}[Bounds for a successful clause parse]
\lean{Std.Sat.CNF.parseClause_bounds}
\label{Std.Sat.CNF.parseClause_bounds}
\leanok
\uses{Std.Sat.CNF.parseClause, Std.Sat.CNF.parseLit, Std.Sat.CNF.parseLit_length}
\difficulty{5}
Let $f \in \mathbb{N}$, let $x$ and $r$ be bit strings and let $C$ be a clause of literals
$(v,b)$ with $v \in \mathbb{N}$ and $b \in \{0,1\}$. Suppose the clause parser with fuel $f$
succeeds on $x$ and returns $C$ with remainder $r$. Then $|r| \le |x|$, and every literal
$(v,b)$ of $C$ satisfies $v + 1 + |r| \le |x|$. Here the clause parser with fuel $f$ on $x$
returns the empty clause and remainder $y$ when $x = 0y$; when $f = g+1$ and
$x = 1^{k+1}0cy$ it prepends the literal $(k,c)$ to the clause returned by the clause parser
with fuel $g$ on $y$, keeping that call's remainder and failing if it fails; and it fails on
every other input.
\end{theorem}

\begin{theorem}[Bounds for a successful formula parse]
\lean{Std.Sat.CNF.parseClauses_bounds}
\label{Std.Sat.CNF.parseClauses_bounds}
\leanok
\uses{Std.Sat.CNF.parseClause, Std.Sat.CNF.parseClause_bounds, Std.Sat.CNF.parseClauses}
\difficulty{5}
Let $f \in \mathbb{N}$, let $x$ and $r$ be bit strings and let $\varphi$ be a CNF formula
over variables indexed by $\mathbb{N}$. Suppose the formula parser with fuel $f$ succeeds on
$x$ and returns $\varphi$ with remainder $r$. Then $|r| \le |x|$, and every literal $(v,b)$
of every clause of $\varphi$ satisfies $v + 1 + |r| \le |x|$. Here the formula parser with
fuel $f$ on $x$ returns the empty formula and remainder $y$ when $x = 0y$; when $f = g+1$
and $x = 1y$ it parses a clause $C$ from $y$ with fuel $g$, leaving $z$, then a formula
$\psi$ from $z$ with fuel $g$, leaving $w$, and returns $C$ followed by the clauses of
$\psi$ with remainder $w$; and it fails otherwise. The clause parser with fuel $g$ on $u$
returns the empty clause and remainder $u'$ when $u = 0u'$, and when $g = h+1$ and
$u = 1^{j+1}0cu'$ it prepends the literal $(j,c)$ to the clause parsed from $u'$ with fuel
$h$, keeping that remainder; it fails otherwise.
\end{theorem}

\begin{theorem}[Variable count from per-literal bounds]
\lean{Std.Sat.CNF.numVars_le_of_literal_bounds}
\label{Std.Sat.CNF.numVars_le_of_literal_bounds}
\leanok
\uses{Std.Sat.CNF.numVars}
\difficulty{4}
Let $\varphi$ be a CNF formula over variables indexed by $\mathbb{N}$ and let
$n \in \mathbb{N}$. If every literal $(v,b)$ of every clause of $\varphi$ satisfies
$v + 1 \le n$, then $N(\varphi) \le n$, where the variable count $N(\varphi)$ is one plus
the largest variable index occurring in a literal of $\varphi$, or $0$ if $\varphi$ contains
no literal.
\end{theorem}

\begin{theorem}[Variable count of a decoded formula]
\lean{Std.Sat.CNF.numVars_decode_le}
\label{Std.Sat.CNF.numVars_decode_le}
\leanok
\uses{Std.Sat.CNF.decode, Std.Sat.CNF.numVars, Std.Sat.CNF.numVars_le_of_literal_bounds, Std.Sat.CNF.parse, Std.Sat.CNF.parseClauses, Std.Sat.CNF.parseClauses_bounds}
\difficulty{4}
For every bit string $x$, the decoded formula satisfies $N(\mathrm{dec}(x)) \le |x|$. Here
the variable count $N(\varphi)$ of a CNF formula $\varphi$ is one plus the largest variable
index occurring in a literal of $\varphi$, or $0$ if no literal occurs, and $\mathrm{dec}(x)$
is the formula obtained by parsing $x$, with recursion fuel $|x|$, according to the grammar
$\langle\text{formula}\rangle ::= 0 \mid 1\,\langle\text{clause}\rangle\,\langle\text{formula}\rangle$,
$\langle\text{clause}\rangle ::= 0 \mid \langle\text{literal}\rangle\,\langle\text{clause}\rangle$,
$\langle\text{literal}\rangle ::= 1^{k+1}\,0\,b$ (the literal with variable index
$k \in \mathbb{N}$ and polarity $b \in \{0,1\}$) when this parse consumes all of $x$, and
the empty formula otherwise.
\end{theorem}
